\scanpage{23}
fibre $(\mathcal X_{\bar s},\underline I_{\bar s})$ est une \inserted{MD-}courbe géométrique sur $k(s)$ -- i.e. $\mathcal X_{\bar s}$ est $0$-connexe, de dim $1$, \deleted{\unsure{ainsi q.}} les pts singuliers sont \inserted{$\notin I_{\bar s}$ et sont} des pts doubles ordinaires, et enfin la condition de rigidité infinitésimale \inserted{(en automorphismes)} est satisfaite -- \deleted{donc} i.e. on exige que $\mathcal X_{\bar s}$ ne soit pas de genre $1$ lisse avec $I_{\bar s}=\varnothing$, \unsure{et} que les comp. irréd. de genre $0$ de $\mathcal X_{\bar s}$ soient de ``poids marqué''
\[
\geq 3.
\]
La différence avec les considérations du n° précédent, c'est que maintenant ($=$ sur une base $\mathcal S$ connexe, p.ex. \deleted{$\unsure{\mathcal S=}$} un trait) les \uline{types} des fibres géométriques peuvent varier d'un pt géom. à un autre. Mais il est clair que le type \uline{numérique} $(g,\nu)$ ne change pas (invariance de $\chi(\widehat X_s,\underline O_{\widehat X_s})\ldots$ et de $\operatorname{card}\underline I_s$), plus précisément,
\margin{\illegible{1}\ \deleted{\illegible{1}}}
\mspage{22}
